\documentclass[a4paper]{article}
\usepackage{amsfonts}
\usepackage{amsmath}
\usepackage{amssymb}
\usepackage{graphicx}     
\usepackage{color}
\usepackage{cancel}


\begin{document}
  
\noindent \large Auteur: Thomas Olszowka \\
\noindent \large Studierichting: Informatica \\
\noindent \large Oefeningenreeks 1D nr. 23.2 \\

\medskip

\normalsize

$iz^3 -(2i - 2)z^2 -3z + i + 1 = 0$\\

\begin{tabular}{c|ccccc}
     & $i$ & $-2i + 2$ & $-3$ & \vline & $i + 1$  \\
     $1$ &  & $i$ & $-i + 2$ & \vline & $-i -1$ \\ \hline
     & $i$ & $-i + 2$ & $-i - 1$ & \vline & $0$ \\
\end{tabular}\\

We bekomen dan de volgende vergelijking $(iz^2 + (-i + 2)z - i - 1)(z-1) = 0$\\

$\Delta = (-i + 2)^2 - 4\cdot i \cdot (-1 -i)$ \\
    
$= -1 - 4i + 4 - 4i(-1 - i)$ \\
    
$= -1 -4i + 4 + 4i + 4i^2$ \\

$= -1 + 4 - 4 - 4i + 4i$ \\
    
$= -1 +\cancel{4} -\cancel{4} -\cancel{4i} +\cancel{4i}$ \\

$= -1$ \\

Dit geeft $(x + yi)^2 = -1$ \\

Vermits $b = 0$ en $a < 0$ is $x = 0$ en $-y^2 = a$. De wortels zijn van de vorm $\left\{i\sqrt{-a}, -i\sqrt{-a} \right\} $, wat ons geeft dat de complexe wortels $\left\{ i, -i \right\}$ zijn.\\

$z_{2,3} = \dfrac{(-2 + i) \pm i}{2i}$\\

$z_2 = \dfrac{-2 +2i}{2i} =  \dfrac{(-2 + 2i)(-2i)}{(2i)(-2i)} =\dfrac{4 + 4i}{4} = 1 + i$ \\

$z_3 = \dfrac{-2}{2i} = \dfrac{(-2)(-2i)}{(2i)(-2i)} = \dfrac{4i}{4} = i$ \\

$z \in \left\{1, 1 + i, i \right\}$ \\

\begin{thebibliography}{999}
\end{thebibliography}
\end{document} 